\section{Mission Overview and Orbit Design}\label{sec:mission-overview}
The ANTARA mission has been conceived as a dedicated scientific CubeSat mission focused on characterizing the high-energy proton radiation environment within the South Atlantic Anomaly (SAA). The spacecraft is intended to operate in a near-circular Sun-Synchronous Low Earth Orbit. The detailed orbit propagation and trade-study results are provided in Appendix~\ref{app:orbit-rf} and Appendix~\ref{app:source-02}.

\subsection{Spacecraft Overview}\label{spacecraft-overview}
ANTARA is a 3U CubeSat with an on-board computer, an electrical power system, an attitude determination and control system, a communications system, and the scientific payload. The initial orbit is approximately 600~km at an inclination of $97.6^\circ$. No propulsion system is planned.

The 3U spacecraft dimensions are approximately $100\times100\times340.5$~mm. The mass budget and 3~kg design limit are assessed in Section~\ref{sec:budget-closure}.

\subsection{Orbit Selection}\label{sec:orbit-design}
The preliminary orbit baseline is a 600~km Sun-Synchronous Orbit with an inclination of $97.6^\circ$ and an eccentricity of approximately 0.001. The latest seven-case RAAN power study reports its largest modeled mean generation at $180^\circ$; power and battery behavior depend on the insertion RAAN. ANTARA does not control the launch RAAN, so the spacecraft is being assessed across multiple RAAN cases. The approved future analysis epoch is 1 January 2029; the evaluated studies retain the source epochs and assumptions identified in Section~\ref{sec:simulation-results}.

\begin{table}[htbp]\centering
\caption{Preliminary orbit parameters}\label{tab:orbit-baseline}
\begin{tabularx}{\linewidth}{@{}lX@{}}\toprule
Parameter & Mission baseline / current status \\ \midrule
Initial altitude & Approximately 600~km \\
Inclination & $97.6^\circ$ \\
Eccentricity & Approximately 0.001 \\
RAAN & Launch-dependent; $0^\circ$ and $30^\circ$ are favorable cases, not requirements \\
Propulsion & None \\
\bottomrule\end{tabularx}\end{table}

The mission targets approximately five years of operational science. The passive-decay study and its propagation endpoint are assessed separately in Section~\ref{end-of-life-deorbit-and-debris-mitigation-compliance}.

\subsection{Mission Lifecycle}\label{sec:mission-lifecycle}
The mission is divided into Launch, Deployment, Cold and Dark, Detumbling, Commissioning, Nominal Operations, Extended Operations, and Mission End phases. Following deployment the spacecraft activates after the planned delay, enters Cold and Dark, performs detumbling, and proceeds toward power-positive operations and commissioning. Nominal operations consist of Sun Pointing, Detection, and Relay, with Safe Mode for recovery.

The extended phase begins after nominal mission objectives have been met. The minimum, full, and extended mission duration targets are defined in Section~\ref{mission-success-criteria}. The detailed phase descriptions and mode transitions are retained in Appendices~\ref{app:source-02} and~\ref{app:operational-logic}.
